\section{Compressing the surfaces with an autoencoder}
\label{sec:autoencoder}

Section~\ref{sec:model} constructed the surface dataset.  Here we transform and
normalize each surface before passing it through three one-dimensional
convolutional layers and a two-number bottleneck.  The decoder reconstructs the
full grid from those two values.  The purpose is not merely to build a neural
network: it is to test whether a nonlinear map can preserve more information
than PCA when both methods are restricted to the same two-number
representation.  Section~\ref{sec:results} makes that matched comparison.

\subsection{Input transformation and normalization}

Large amplitudes would dominate an error measured directly on the positive
surface, even when a smaller-amplitude region has an equally important
fractional error.  The encoder therefore receives the pointwise log-surface
$Y=\log_{10}\mathcal A_\Theta$, which turns multiplicative differences into
additive ones.  We then subtract the mean
training surface $\overline{Y}$ and divide by its global root-mean-square
deviation $\sigma$.  Both quantities are computed from the training split and
then applied unchanged to validation and test data.  Using one global scale,
rather than a different scale at each $(k,z)$, keeps the relative size of
features across the grid and avoids leakage from the held-out subsets.  The
full transformation is recorded in Appendix~\ref{app:detailed-equations}.

\subsection{From 3,131 values to a two-number bottleneck}

The bottleneck is the part that tests the project's compression question.
An autoencoder contains an \emph{encoder} and a \emph{decoder}.  The encoder
maps a surface $\bm{x}$ to a latent vector $\bm{z}$, and the decoder maps
$\bm{z}$ to a reconstruction $\widehat{\bm{x}}$.  We set
$\bm{z}\in\mathbb{R}^{2}$ to match the two-component PCA baseline and to make
this pilot representation easy to inspect.  Because no skip connection
bypasses the bottleneck, every reconstructed surface must genuinely pass
through two numbers.  Otherwise good reconstruction could come from a hidden
high-dimensional route rather than compression.  The compact mathematical
definition is included in Appendix~\ref{app:detailed-equations}.

One-dimensional convolutions first collect nearby structure along wavenumber.
Dense layers then summarize the whole surface.  The bottleneck keeps only two
coordinates, from which the decoder must rebuild the full grid.  Successful
reconstruction would therefore indicate that the sampled surfaces lie near a
lower-dimensional nonlinear structure; failure would show that two
coordinates discard too much information.

\subsection{Why use one-dimensional convolutions?}

Neighboring wavenumbers are strongly related because the analytic response
changes smoothly and contains scale transitions.  We therefore arrange the
input as 31 redshift channels, each containing a 101-point sequence in
wavenumber.  A one-dimensional convolution slides along that sequence and
learns local scale-dependent patterns.  Reusing the same kernel at different
$k$ positions lets the network recognize similar transitions without learning
a separate weight for every grid cell, while its filters mix information
between redshift channels.  This is a motivated first architecture, not a
claim that Conv1D is uniquely optimal; its advantage must be judged against
the held-out reconstruction results.

\subsection{Network architecture}

\begin{table}[t]
\centering
\caption{Configured three-layer Conv1D autoencoder with latent dimension 2.
Sequence lengths refer to the
wavenumber axis.}
\label{tab:architecture}
\begin{tabular}{@{}lll@{}}
\toprule
Stage & Shape or width & Main operation \\
\midrule
Input & $31\times101$ & normalized $\log_{10}\mathcal A_\Theta$ \\
Encoder 1 & $64\times51$ & Conv1D, kernel 5, stride 2, SiLU \\
Encoder 2 & $128\times26$ & Conv1D, kernel 5, stride 2, SiLU \\
Encoder 3 & $256\times13$ & Conv1D, kernel 3, stride 2, SiLU \\
Flatten & $3{,}328$ & flatten $256\times13$ features \\
Dense encoder & $3{,}328\rightarrow256\rightarrow64$ & nonlinear summary \\
Bottleneck & $2$ & latent coordinates \\
Dense decoder & $2\rightarrow64\rightarrow256\rightarrow3{,}328$ & expand latent representation \\
Decoder & reverse sequence & transposed convolutions \\
Output & $31\times101$ & reconstructed normalized surface \\
\bottomrule
\end{tabular}
\end{table}

The encoder gradually shortens the wavenumber sequence while increasing the
number of learned features; the decoder reverses this process.  This gives the
network enough capacity to test a nonlinear representation without relaxing
the two-number information bottleneck.  The resulting model has 2,039,969
trainable weights.  That large weight count should not be confused with the
latent dimension: the project tests compact representation of each surface,
not a network with only two trainable parameters.  We use SiLU activations and
no skip connections; a skip connection would let information avoid the
bottleneck and invalidate the comparison with two-component PCA.

\subsection{Training objective and model selection}

The training target should reward accurate reconstruction over the whole
surface rather than at one selected scale.  We therefore minimize mean squared
error in normalized log-amplitude space, averaged over every sample, redshift,
and wavenumber value in a batch.  Its explicit definition is given in
Appendix~\ref{app:detailed-equations}.  Table~\ref{tab:training-config}
shows only the choices needed to understand the training logic; complete
numerical settings are listed in Appendix
Table~\ref{tab:training-config-details}.
These settings describe the current shared template.  The preliminary
autoencoder metrics in Section~\ref{sec:results} come from a separate
25-epoch run, not from a completed run of this longer schedule; the table must
not be read as the provenance of those reported numbers.

\begin{table}[htbp]
\centering
\small
\caption{Main choices in the shared autoencoder training template.}
\label{tab:training-config}
\begin{tabularx}{\linewidth}{@{}
  >{\raggedright\arraybackslash}p{0.19\linewidth}
  >{\raggedright\arraybackslash}p{0.25\linewidth}
  >{\raggedright\arraybackslash}X@{}}
\toprule
Setting & Configured value & Purpose \\
\midrule
Objective & MSE in normalized $\log_{10}\mathcal A_\Theta$
  & Reconstruct the full surface while treating multiplicative amplitude
    differences more evenly. \\
Optimizer & AdamW, learning rate $2\times10^{-4}$
  & Learn nonlinear weights with adaptive updates. \\
Batch size & 512 surfaces
  & Balance stable updates with memory use. \\
Training control & Reduce the learning rate after a plateau; stop early
  & Refine stalled optimization without training indefinitely. \\
Model selection & Lowest validation MSE
  & Select performance on unseen surfaces rather than training fit. \\
Success benchmark & Recovered variance $\geq 99.9\%$
  & Set a demanding compression goal, not a cosmological requirement. \\
\bottomrule
\end{tabularx}
\end{table}
\clearpage

\subsection{What must be checked beyond the average loss?}

An average MSE can look good even if the model fails badly for a small region
or a difficult surface.  Such a failure matters because a localized IA error
could propagate differently into a lensing analysis.  We therefore use
several complementary diagnostics:

\begin{description}
  \item[Log-space RMSE] measures the typical reconstruction residual
  $\widehat Y-Y$ in dex over all validation grid values.  It is closely tied
  to the training objective, but its logarithmic units are not an immediate
  fractional error in physical amplitude.

  \item[Recovered variance] is one minus the reconstruction sum of squares
  divided by the validation set's total squared distance from the training
  mean.  It answers how much of the surface-to-surface variation survives
  compression and permits a direct comparison with PCA.  A high value can
  still conceal a localized failure.

  \item[Mean physical relative error] converts each log residual back to
  amplitude space using
  $\lvert 10^{(\widehat Y-Y)}-1\rvert$ and averages over surfaces and grid
  cells.  This quantity is easier to interpret fractionally, but positive and
  negative log residuals produce asymmetric amplitude errors.

  \item[Per-surface relative RMSE] first summarizes the physical relative
  errors across all 3,131 cells of each surface.  Its 95th percentile then
  describes a difficult but not single worst-case part of the validation
  population.  This prevents a small set of poorly reconstructed surfaces
  from disappearing inside a global mean.
\end{description}

Error maps over $(k,z)$ and parameter space are also needed to locate where a
model fails, while maximum-error summaries can flag isolated extreme cells.
The preliminary run available for this report provides only the scalar
summaries reported in Section~\ref{sec:results}; its checkpoint is unavailable
for generating new maps.  Finally, the $99.9\%$ recovered-variance target is
an internal compression benchmark, not evidence that reconstruction errors
are already small enough for cosmology.  That requires propagating them into
lensing observables and parameter inference.
